\documentclass[10pt]{article}
\usepackage{amsmath}
\usepackage{amsthm}
\usepackage{amssymb}
\usepackage{fullpage}
%\usepackage{enumerate}
\usepackage{graphicx}
\graphicspath{ {./} }
\usepackage{hyperref}
\usepackage{multirow}
%\usepackage[]{algorithm2e}
\usepackage{minted}
\usepackage{tikz}

\usetikzlibrary{arrows,automata,positioning}

\newcommand{\abs}[1]{{\lvert #1\rvert}}
\newenvironment{solution}{\textbf{Solution:}}{\hfill$\square$}
%\newcommand{\pr}{\textsc{P}}

\begin{document}

\begin{flushright}
\begin{tabular*}{\textwidth}
{@{\extracolsep{\fill}}l r}
% \hline
\multirow{3}{0.7\textwidth}{{\Huge CS687A: Assignment 1}} & \\ 
& \\ 
& \large{Kumar Arpit: 200532}\\
% & Kumar Arpit: 200532\\
\hline
\end{tabular*}
\end{flushright}

\begin{enumerate}

    \item Prove that every infinite computably enumerable language contains an infinite decidable subset.

    \begin{solution}

    To prove that every infinite computably enumerable (c.e.) language contains an infinite decidable subset, we can use the fact that a language is c.e. if and only if it is the range of a total computable function. Let $L$ be an infinite c.e. language.
    
    We then construct a decider for a subset of $L$ using dovetailing. Let $M$ be the Turing machine that enumerates the language $L$. We can simulate $M$ on all possible inputs in parallel by interleaving its computation steps, in a way that ensures that each input is eventually simulated infinitely often as follows

    \begin{enumerate}
        \item Start by simulating $M$ on input 1 for one step.
        \item Then simulate $M$ on input 1 and input 2 for one step each.
        \item Then simulate $M$ on input 1, 2, and 3 for one step each.
        \item Continue this way, simulating $M$ on inputs 1 to $k$ for one step each at step $k$, for all $k > 0$. By this procedure, we simulate $M$ on all inputs infinitely often, and we can construct a decider for a subset of $L$ as follows:
    \end{enumerate}

    At step $k$ of the dovetailing, check if $M$ has output a string in $L$ on any input in $\{1,2,\dots,k\}$. If it has, output that string as a member of the decidable subset of $L$.
    
    If $M$ has not output a string in $L$ on any input in $\{1,2,\dots,k\}$, continue to the next step of the dovetailing.
    
    The resulting subset of $L$ is clearly decidable, since we can decide whether any string is a member of the subset by running the above algorithm for long enough to see whether that string is output.
    
    To see that the resulting subset is infinite, note that since $L$ is infinite and $M$ eventually simulates every input infinitely often, there must be infinitely many inputs on which $M$ outputs a string in $L$. Therefore, the decidable subset we construct contains infinitely many strings.
        
    \end{solution}
    \newpage

    \item  Let $x_0x_1\ldots x_{n-1}$ be a finite binary string with $n$ bits. For $1\leq k\leq n - 1$, let the “$k^{th}$-left rotation of $x$” be $x_k \ldots x_{n-1}x_0\ldots x_{k-1}$ (corresponding to left-shift and rotate on $x$ done $k -1$ times). Do all rotations of x have the same plain Kolmogorov complexity of x? If so, prove your claim. Otherwise, construct a counterexample, and prove that your string has some left rotation which has a significantly different Kolmogorov complexity.

    \begin{solution}

    To prove that all rotations of $x$ have the same plain Kolmogorov complexity as $x$, we need to show that given any rotation $y$ of $x$, we can construct a program of length at most $c + O(1)$ bits that outputs $y$, where $c$ is the plain Kolmogorov complexity of $x$.
    
    To do this, let $y$ be the $k$-th left rotation of $x$. We can construct a program $P$ that outputs $y$ as follows:

    \begin{enumerate}
        \item Output the length $n$ of $x$.
        \item Output the integer $k$.
        \item Output the binary string $x$ (using a short program for it).
        \item In our program, we also have a function that takes as input an integer $k$ and a binary string $s$, and outputs the $k$-th left rotation of $s$. This function is implemented using a loop that iteratively shifts and rotates the string $s$ $k$ times i.e. it is of constant length.
        \item Apply the function from step (d) to the input $(k, x)$.
    \end{enumerate}
    
    The length of this program (say $d$) is $d \leq c + O(\log(k)) + O(1)$, where $c$ is the plain Kolmogorov complexity of $x$. With a parallel argument, it's  easy to see $c \leq d + \log(n-k) + O(1)$ where $d$ is the plain Kolmogorov complexity of $y$. Therefore, we have shown that given any rotation $y$ of $x$, we can construct a program($d$) of length at most $d = c + O(1)$ bits that outputs $y$ (and vice versa), which implies that all rotations of $x$ have the same plain Kolmogorov complexity as $x$.
        
    \end{solution}
    \newpage

    \item Let $X_H$ be the infinite binary sequence defined as follows. For any number $i$, the $i^{th}$ bit of $X_H$ is 1 if the $i^{th}$ Turing machine halts, and 0 otherwise. Since the Halting problem is uncomputable, $X_H$ is uncomputable. However, show that prefixes of $X_H$ are compressible: for all sufficiently large n, show that $C(X_H[0 \ldots n - 1]) \leq \log n + O(1)$, where $X[0 . . . n - 1]$ denotes the $n$-length prefix of $X_H$. (This shows that even when a string is uncomputable, it may be highly compressible.)

    \begin{solution}

    We will construct a program that compresses the prefixes of $X_H$.
    
    Let $n$ be the length of the prefix $X_H[0 \ldots n - 1]$ and let $M_1, M_2, \ldots, M_k$ be an enumeration of all Turing machines with $k$ states. For each $i \leq k$, we define $p_i$ to be the maximum number of steps that $M_i$ takes to halt on any input of length at most $n$ (if $M_i$ does not halt on any such input, then we set $p_i$ to be $n$). Note that each $p_i$ is at most exponential in $k$.
    
    Our compression program works as follows:
    
    On input $X_H[0 \ldots n - 1]$, it outputs the index $i$ of the first Turing machine $M_i$ such that $M_i$ halts on an input of length at most $n$ and runs for exactly $p_i$ steps before halting. This index can be represented using $\log k$ bits. Since $k$ is exponential in $p_i$, we have $\log k = O(\log p_i)$. Since each $p_i$ is at most exponential in $k$, we have $\log p_i = O(\log k)$. Therefore, the size of the output of our program is $O(\log k) = O(\log p_i)$.
    
    We now show that our program compresses prefixes of $X_H$.
    
    Let $X_H[0 \ldots n - 1]$ be a prefix of $X_H$ and let $i$ be the index output by our program. We claim that $M_i$ halts on an input of length at most $n$ and runs for exactly $p_i$ steps before halting. Suppose, for the sake of contradiction, that $M_i$ halts on an input of length at most $n$ but does not run for exactly $p_i$ steps before halting. Then there exists an input $x$ of length at most $n$ on which $M_i$ halts and runs for fewer than $p_i$ steps, or $M_i$ does not halt on any input of length at most $n$.
    
    If $M_i$ halts on $x$ and runs for fewer than $p_i$ steps, then the output of $M_i$ on $x$ can be computed in fewer than $p_i$ steps, which contradicts the definition of $p_i$. Therefore, $M_i$ does not halt on any input of length at most $n$. But then $X_H[0 \ldots n - 1]$ does not contain the $i^{th}$ bit of $X_H$, which is 1 if and only if $M_i$ halts on some input. This contradicts the fact that $i$ was output by our program.
    
    Therefore, our program correctly outputs the index $i$ of a Turing machine that halts on an input of length at most $n$ and runs for exactly $p_i$ steps before halting. The size of the output is $O(\log p_i)$. Since each $p_i$ is at most exponential in $k$, the size of the output is $O(\log k)$.
    
    Thus, the prefix $X_H[0 \ldots n - 1]$ is compressible to $O(\log n + O(1))$ bits, as required.


        
    \end{solution}
    \newpage

    \item Let $\omega$ be an infinite binary sequence defined by
    \[\omega=\sum_{i\in \mathbb{N}, M_i(i)\downarrow}\frac{1}{2_i}\]
    Show that prefixes of $\omega$ are incompressible: for all sufficiently large $n$, $C(\omega[0 \ldots n - 1]) \geq n-O(1)$. (The purpose of this question is to show that there are uncomputable strings which are also incompressible, in contrast to the previous question.)

    \begin{solution}

    Let $s=\omega[0 \ldots n-1]$ be a prefix of $\omega$ of length $n$. We will show that $C(s) \geq n-O(1)$.
    
    First, note that $s$ can be written as a sum of the form
    
    \[s = \sum_{i=1}^n \frac{a_i}{2^i}\]
    
    where $a_i \in \{0,1\}$ are the individual bits of $s$. Since $\omega$ is defined as the sum of all finite binary strings that correspond to halting Turing machines, we know that there exists a Turing machine $M$ that halts and outputs $s$.
    
    To encode $s$, we can encode the binary representation of each $a_i$ as a string of $k = \lceil \log_2 n \rceil$ bits. We can then concatenate these $k$-bit strings together to form a single string $x$ of length $nk$, which is a binary encoding of $s$.
    
    Therefore, we have $|x| = nk$ and $s = \text{decode}(x)$, where $\text{decode}(x)$ is the function that maps a binary string $x$ of length $nk$ to the corresponding sum $\sum_{i=1}^n \frac{a_i}{2^i}$.
    
    To show that $C(s) \geq n-O(1)$, it is sufficient to give a string $y$ such that $|y| \leq n-O(1)$ and $y$ is not prefix-free with respect to $x$. If such a string exists, then $x$ cannot be compressed to a length much less than $n$, since the compressor would have to include $y$ in its output.
    
    Consider the following string $y$:
    
    \[y = 1^{k}0^{k-1}1^{k}0^{k-2}1^{k}\ldots 0^{1}1^{k}\]
    
    where there are $n$ stretches of $k$ bits each, and the $i$th stretch contains $k-i+1$ ones followed by $i-1$ zeros. In other words, the $i$th stretch contains the binary representation of the number $2^i-1$.
    
    Note that $|y| = nk - \frac{n(n+1)}{2}$, which is approximately equal to $n^2/2$.
    
    We claim that $y$ is not prefix-free with respect to $x$. Indeed, consider the $i$th block of $y$, which corresponds to the $i$th bit of $s$. This block consists of $k-i+1$ ones followed by $i-1$ zeros, which is the binary representation of $2^i-1$. Thus, the $i$th block of $y$ is also a valid encoding of the number $2^i-1$ in binary.
    
    However, the binary representation of $2^i-1$ can also be obtained by concatenating the first $i$ bits of the binary representation of $s$. Therefore, the $i$th block of $y$ is a prefix of the $i$th block of $x$, which is the $k$-bit encoding of $a_i$.
    
    Since this holds for all $i \in {1, 2, \ldots, n}$, we have that $y$ is not prefix-free with respect to $x$. Therefore, $C(x) \geq |x| - |y| \geq n - O(1)$, which implies that $C(s) \geq n - O(1)$, as desired.
    
    Therefore, we have shown that for any sufficiently large $n$, the Kolmogorov complexity of the first $n$ bits of $\omega$ is at least $n-O(1)$. Hence, prefixes of $\omega$ are incompressible.
        
    \end{solution}
    \newpage

    \item Let $f : \Sigma^*\rightarrow \Sigma^*$ be a total computable function. Then show that for every string $x$, we have $C(f(x)) \leq C(x) + O(1)$. (The purpose of this question is to show that you can never substantially increase the information in any string through computation: you can only destroy the information.)

    \begin{solution}

    We can prove this using the concept of Kolmogorov complexity. Let $x$ be any string in $\Sigma^*$, and let $y = f(x)$.
    
    We know that the Kolmogorov complexity of a string $s$, $C(s)$, is the length of the shortest program that outputs $s$ when run on a universal Turing machine.
    
    Let $M$ be a Turing machine that computes the function $f$, and let $p$ be the program that represents $M$. Then we can construct a program $q$ that takes $x$ as input, runs $p$ on $x$, and outputs the result $y$. The length of $q$ is $|q| = |p| + |x| + O(1)$, since $q$ consists of the program $p$, the input $x$, and some fixed overhead.

Thus, we have:

$$C(y) \leq |q| \leq |p| + |x| + O(1) = C(M) + |x| + O(1)$$

Since $M$ is a fixed program that computes $f$ for all inputs, its complexity is also fixed i.e. a constant. So we can absorb $C(M)$ into the $O(1)$ term. Therefore, we have:

$$C(y) \leq C(x) + O(1)$$

This holds for any string $x$, so the statement is true for all inputs.


        
    \end{solution}
    \newpage

    \item Show, using an argument using plain Kolmogorov complexity, that for all sufficiently large $n$, most strings of length $n$ will not have a contiguous stretch of more than $\log_2 n$ zeroes anywhere in the string. (This shows how you can prove that some event happens with very high probability, using the theory of plain Kolmogorov complexity.)

    \begin{solution}

    We will prove the statement by contradiction, i.e., we will assume that there exists a sufficiently large $n$ such that more than half of the strings of length $n$ contain a contiguous stretch of $\log_2 n$ or more zeroes, and show that this leads to a contradiction with the definition of plain Kolmogorov complexity.

Let $S$ be the set of all binary strings of length $n$, and let $S'$ be the subset of $S$ that contains all strings with a contiguous stretch of $\log_2 n$ or more zeroes. By assumption, $|S'| > \frac{1}{2}|S|$.

Consider the function $f$ that maps each string $x$ in $S$ to the shortest program (in some fixed programming language) that outputs $x$. By the definition of plain Kolmogorov complexity, the Kolmogorov complexity $K(x)$ of $x$ is equal to the length of $f(x)$, plus a constant that depends only on the programming language. Therefore, $K(x) \leq |f(x)| + c$ for some constant $c$.

Now, consider the set $T$ of all programs $p$ that output a string in $S'$. Since $|S'| > \frac{1}{2}|S|$, it follows that there are more than $2^{n-1}$ strings in $S'$, and therefore there are more than $2^{n-1}$ distinct programs in $T$. Each program in $T$ has length at most $\log_2 n$ more than the length of the string it outputs, since any program that outputs a string with a contiguous stretch of $\log_2 n$ or more zeroes can simply output the zeroes explicitly. Therefore, the total length of all programs in $T$ is at most $2^{n-1}(\log_2 n + c)$.

Since $T$ is a set of programs that output strings in $S'$, and $|S'| > \frac{1}{2}|S|$, it follows that there are more strings in $S'$ than there are distinct programs in $T$. Therefore, by the pigeonhole principle, there must be at least one string $y$ in $S'$ that can be output by more than one program in $T$. Let $p_1$ and $p_2$ be two such programs that output $y$.

Then, we can concatenate the first $\log_2 n$ bits of $p_1$ with the last $|p_2| - \log_2 n$ bits of $p_2$ to form a new program $p'$ that outputs $y$. The length of $p'$ is at most $\log_2 n + |p_2| - \log_2 n = |p_2|$. Therefore, we have $K(y) \leq |p'| + c \leq |p_2| + c$.

However, $p_1$ and $p_2$ are distinct programs that output $y$, so $|p_1| \neq |p_2|$. Without loss of generality, assume that $|p_1| < |p_2|$. Then, we have $K(y) \leq |p_2| + c < 2^{n-1}(\log_2 n + c) + c = O(n\log n)$.

On the other hand, since $y$ has a contiguous stretch of $\log_2 n$ or more zeroes, we have $K(y) \geq K(\text{stretch of zeroes}) + c$. Let $z$ be the shortest program that outputs the stretch of zeroes in $y$. Then, we have $K(y) \geq |z| + c$, where $|z|$ is at most $\log_2 n$. Therefore, we have $|z| + c \leq K(y) \leq |p_2| + c$.

This leads to a contradiction, since the length of the stretch of zeroes in $y$ is at least $\log_2 n$, which means that $|z| \geq \log_2 n$. Therefore, we have $\log_2 n + c \leq K(y) < O(n\log n)$, which implies that $n$ can be recovered from $y$ with high precision using a short program. This contradicts the assumption that more than half of the strings of length $n$ contain a contiguous stretch of $\log_2 n$ or more zeroes, since it implies that $n$ has low Kolmogorov complexity relative to the set of all strings of length $n$.

Therefore, we conclude that for all sufficiently large $n$, most strings of length $n$ will not have a contiguous stretch of more than $\log_2 n$ zeroes anywhere in the string, since otherwise it would lead to a contradiction with the definition of plain Kolmogorov complexity.


        
    \end{solution}
    \newpage

    \item  Show that $C(xy) \neq C(x) + C(y) + O(1)$.

    \begin{solution}

        Consider the two languages $L_1=01^n$ and $L_2=1^m$.
        
        The complexities of $L_1$ and $L_2$ are bounded by $C_1=\log(n)$ and $C_2=\log(m)$.

        However, to differentiate the concatenation of two strings in $L_1$ and $L_2$ we require at least $\log(\log(n))$ bits for any possible encoding scheme.

        Thus the plain complexity of the concatenation of the two strings is strictly greater than the sum of their individual complexities.

        That is we have,

        \[C(xy)>C(x)+C(y)+O(1)\]
        
    \end{solution}
    \newpage
    
\end{enumerate}

\end{document}